\section{Results}

\begin{frame}{Accuracy holds up until 50\% noise}
  \begin{columns}[T]
    \begin{column}{0.55\textwidth}
      \centering
      \includegraphics[width=\linewidth,height=0.65\textheight,keepaspectratio]{acc_vs_noise.png}
    \end{column}
    \begin{column}{0.45\textwidth}
      \small
      \begin{itemize}
        \item Stays flat up to $10\%$ noise ($0.961 \to 0.934$)
        \item Drops sharply at $50\%$, down to $0.708$
        \item The model can handle about 1 bad verb in 10 before it breaks
        \item Not a gradual decline, more like a threshold
      \end{itemize}
    \end{column}
  \end{columns}
\end{frame}

\begin{frame}{The attractor gap grows with noise}
  \begin{columns}[T]
    \begin{column}{0.55\textwidth}
      \centering
      \includegraphics[width=\linewidth,height=0.65\textheight,keepaspectratio]{attractor_gap.png}
    \end{column}
    \begin{column}{0.45\textwidth}
      \small
      \begin{itemize}
        \item Gap $=$ acc(no attractor) $-$ acc(attractor)
        \item Goes from $0.073$ to $0.198$ as noise increases
        \item Tells us how much the model leans on the nearest-noun shortcut
        \item More noise means more reliance on the linear rule
      \end{itemize}
    \end{column}
  \end{columns}
\end{frame}

\begin{frame}{Attractor items are where it falls apart}
  \begin{columns}[T]
    \begin{column}{0.55\textwidth}
      \centering
      \includegraphics[width=\linewidth,height=0.65\textheight,keepaspectratio]{acc_by_attractor.png}
    \end{column}
    \begin{column}{0.45\textwidth}
      \small
      \begin{itemize}
        \item Without attractors, accuracy stays above $0.94$ until $50\%$
        \item With attractors, it drops to near chance ($0.525$)
        \item Attractor items are the only cases where the two rules disagree
        \item So that's exactly where the model struggles first
      \end{itemize}
    \end{column}
  \end{columns}
\end{frame}

\begin{frame}{Learning over time: delay vs.\ prevention}
  \begin{columns}[T]
    \begin{column}{0.55\textwidth}
      \centering
      \includegraphics[width=\linewidth,height=0.65\textheight,keepaspectratio]{acc_trajectory.png}
    \end{column}
    \begin{column}{0.45\textwidth}
      \small
      \begin{itemize}
        \item Up to $10\%$ noise: reaches the same accuracy, just takes longer (\textbf{delay})
        \item At $50\%$: plateaus around $0.69$ and never catches up (\textbf{prevention})
        \item Clean model: gap shrinks over training ($0.15 \to 0.05$)
        \item $50\%$ model: gap actually \emph{grows} ($0.14 \to 0.20$)
      \end{itemize}
    \end{column}
  \end{columns}
\end{frame}

\begin{frame}{Does structure help? RNNG vs.\ LSTM}
  \begin{columns}[T]
    \begin{column}{0.55\textwidth}
      \centering
      \includegraphics[width=\linewidth,height=0.65\textheight,keepaspectratio]{combined_acc_vs_noise.png}
    \end{column}
    \begin{column}{0.45\textwidth}
      \small
      \begin{itemize}
        \item RNNG $=$ LSTM plus an explicit syntactic (tree) bias
        \item More accurate than the LSTM at \emph{every} noise level
        \item Gap widens where noise bites hardest: at $50\%$, $0.834$ vs.\ $0.708$
        \item Same held-out set ($\sim$20k pairs); single full-scale seed so far
      \end{itemize}
    \end{column}
  \end{columns}
\end{frame}

\begin{frame}{RNNG keeps the hierarchical rule under noise}
  \begin{columns}[T]
    \begin{column}{0.55\textwidth}
      \centering
      \includegraphics[width=\linewidth,height=0.65\textheight,keepaspectratio]{combined_attractor_gap.png}
    \end{column}
    \begin{column}{0.45\textwidth}
      \small
      \begin{itemize}
        \item LSTM's attractor gap grows $2.7\times$ with noise ($0.073 \to 0.198$)
        \item RNNG's stays low and flat ($0.020 \to 0.036$)
        \item At $50\%$ noise the RNNG gap is \emph{half} the LSTM's clean-data gap
        \item Structural bias preserves the hierarchical rule where the LSTM defaults to the linear shortcut
      \end{itemize}
    \end{column}
  \end{columns}
\end{frame}
